% Diese Fan-Geschichte ist noch in Arbeit. Ich teile sie bereits hier, damit sie für mich nicht verloren geht, falls mein Laptop spontan verschwinden sollte, und damit sie für euch nicht völlig verloren geht, falls ich spontan verschwinden sollte. Fühlt euch in letzterem Fall vollkommen frei, hierrin zu schmökern und nach Lust und Laune weiter daran zu werkeln. Aber solange ich ein aktiver Andori bleibe, lest bitte nicht zu genau nach, was ich geplant haben könnte. LG BBB :D

\begin{chapterbox}
    \chapter{Die Falschen Drachen (WIP)}
    \label{wip-falschen-drachen}
    \az{76 bis 77}

    Der Tross der Andori erreicht die Ödnis Krahd, um die dort versklavten Ambacus und entführten Andori zu befreien. Die Gefangenschaft in der Winterburg hat Iolith zum Drachenkult konvertiert. Die Überbleibsel der Crew der Schwarzen Kogge haben sich übers Andorversum verteilt, manche von ihnen auf Rache gegenüber anderen sinnend. Mart und Bort erhalten Besuch von \gend{Hüter*innen}{Hütern} der Zeit. 
\end{chapterbox}





\section{Borghorns Fall}
\az{76 bis 77}

Griun erlebte den Angriff des Tross der Andori auf Borghorn mit und konnte im letzten Moment fliehen.






\newpage
\section{Tode und Tote}
\az{77}

Chada gedenkte den im Kampf Gefallenen.







\newpage
\section{Bort}


Zwei \gend{Hüter*innen}{Hüter} der Zeit besuchten Cavern.







\newpage
\section{Mart}

Ein Zeitportal aus der tiefsten Vergangenheit des Zwergenreichs öffnete sich.








\newpage
\section{Thogger}

Meres besuchte Thogger, um ihn davon zu überzeugen, Callem aufzuspüren.









\newpage
\section{Meres}


Meres und Thogger stellten sich Callem.






\newpage
\section{Kardòmal}



Kardòmals Partnerin Griun kämpfte mit ihrer neuen Identität. Die Stimmen der Drachen wurden lauter.












\newpage
\section{Callem}


Callem floh vor seinen Untaten. Die völlig aufglöste Kardòmal zu treffen, stellte ihn vor eine fundamentale Entscheidung.











\newpage
\section{Griun}

\az{76}

Griun war zurück in ihrem schlimmsten Albtraum. Doch sie konnte Krahd auch ein zweites Mal entkommen. Und zurück zu ihrer Partnerin finden.














\newpage
\section{Naraven}

Besuch traf ein bei Naravens Hütte. Sein Hadrischer Spiegel sollte Iolith, Griun, Fir und Callem erlauben, in die Vergangenheit zu blicken. 

Callem wusste natürlich noch, wo alle Verstecke der Schwarzen Kogge lagen. Doch dasjenige, das sie betreten wollten, konnte er nicht ohne Weiteres aufsuchen, nicht mehr seit der Befreiung Narkons. Weil ein gewisser anderer Pirat das Versteck mit allerlei trickreichen Fallen geschützt hatte.
Nabib hatte Jari Dorrs Schatzkarte gesehen, vor Jahrzehnten, als der Handelszwerg Garz sie ihm gezeigt hatte. Er könnte sich heute ebenso wenig an die Details erinnern wie Garz, aber seine Knochen durchaus.














\newpage
\section{Fir}

Die Gruppe reiste weiter zum gewünschten Versteck.

Iril und andere Heldinnen kamen ihnen zuvor.





\newpage
\section{Iril}


Iril hatte endlich die \gend{Kultist*innen}{Kultisten} aufgespürt, welche ihre Schwester entführt hatten. Doch ihre Schwester war gar nicht unglücklich mit ihrer Gesellschaft. Iril war hingegen überhaupt nicht glücklich mit der Anwesenheit des Mörders Callem.








\newpage
\section{Die Verfluchte}


Die verfluchte Ambacu eroberte die Rietburg.








\newpage
\section{Iolith}


Die Falschen Drachen boten Iolith ungeahnte Macht und teleportierten sie dafür nach Krahd.








\newpage
\section{Nalledor}

Nalledor folgte den Spuren der Spiegelerinnerungen ins zukünftige Krahd.









\newpage
\section{Epiloge}















\newpage
\section{Appendix: Eine Chronik eines Nekromanten}






\newpage
\section{Appendix: Eine Chronik der Falschen Drachen}


